\chapter{Gradient-Based Samplers}
\label{ch:gradient}

\newthought{Gradient-based samplers use the slope} of the log-likelihood to take long, informed
steps, which makes them very efficient on smooth, high-dimensional posteriors. \Jarvis\ defines
three---\sampler{MALA}, \sampler{HMC}, and \sampler{NUTS}---and they accept the standard
\textsc{mcmc} \yamlkey{Bounds}.

\begin{jwarn}
\textbf{Status: placeholder.} In the current version the gradient \emph{engines} for
\sampler{MALA}, \sampler{HMC}, and \sampler{NUTS} are placeholders---the likelihood-gradient
contract is not yet wired up. Treat this chapter as documentation of the intended methods and their
configuration, not as production-ready samplers. For posterior sampling today, use the
gradient-free \textsc{mcmc} methods in Chapters~\ref{ch:mcmc-1}--\ref{ch:mcmc-3}.
\end{jwarn}

\section{The shared idea}
\label{sec:grad-idea}

A random walk wanders; a gradient walk \emph{climbs}. By following the gradient of the
log-likelihood, these methods propose points that are both far from the current one and likely to
be accepted, so they decorrelate quickly and scale well with dimension---provided a reliable
gradient is available. They use the same \yamlkey{Bounds} as the rest of the family
(\yamlkey{num\_chains}, \yamlkey{num\_iters}, \yamlkey{proposal\_scale}, the last setting the
step/leapfrog size).

\section{MALA}
\label{sec:s-mala}

\paragraph{How it works} The Metropolis-adjusted Langevin algorithm proposes a step that combines a
nudge \emph{up the gradient} with random noise (a discretised Langevin diffusion), then applies a
Metropolis accept/reject to correct the discretisation. It is the simplest gradient method---one
gradient evaluation per step.

\section{HMC}
\label{sec:s-hmc}

\paragraph{How it works} Hamiltonian Monte Carlo gives the point a random ``momentum'' and then
simulates frictionless physics---rolling along the posterior surface using the gradient via several
leapfrog steps---to propose a distant point, which is accepted with a Hamiltonian Metropolis rule.
Long trajectories mean nearly independent samples, but the trajectory length and step size need
tuning.

\section{NUTS}
\label{sec:s-nuts}

\paragraph{How it works} The No-U-Turn Sampler\cite{Hoffman:2011ukg} is \sampler{HMC} that tunes
itself: it keeps extending the trajectory until the path starts to double back (a ``U-turn''), so
you never have to choose the number of leapfrog steps. It is the usual default among gradient
samplers when gradients are available.

\section{Configuration}
\label{sec:grad-format}

When the engines are completed, all three will take the standard form:

\begin{yamlwin}
Sampling:
  Method: "NUTS"          # or MALA, HMC
  Variables: [ ... ]
  Bounds:
    num_chains: 4
    num_iters: 1000
    proposal_scale: 0.1   # step / leapfrog size
  LogLikelihood:
    - {name: "LogL", expression: "..."}
\end{yamlwin}
